\documentclass[11pt,a4paper]{coverletter}

\senderblock
  {Radheem Bin Razi}
  {Distributed Systems \& Cloud-Native Engineer}
  {Ilmenau, Germany \textbar\ \href{mailto:sheikh.radheem@gmail.com}{sheikh.radheem@gmail.com} \textbar\ \href{https://radheem.github.io/my-cv}{portfolio}}

\begin{document}

%% ===== ENGLISH =====
\recipient{XpertDirect}{Hiring Team}{}

\opening{Dear Hiring Team,}

Building backend services that stay fast and correct under heavy transaction volume is the work I know best. At Bluefin Exchange I was the design authority for the core services of a high-throughput crypto exchange — a domain where every API call carries money, latency is unforgiving, and reliability is non-negotiable. That is the same shape of problem behind a high-volume digital commerce platform serving customers across Europe, and it is exactly the kind of role described here.

For just over two years at Bluefin I owned core backend services built on NodeJS, TypeScript, Go and PostgreSQL. I drove system-design improvements independently, reviewed the team's code and design decisions, and kept release processes and documentation in good order while mentoring junior engineers. The platform ran on Docker and Kubernetes with Kafka handling event-driven messaging between services — directly relevant to your nice-to-haves around event-driven architectures and microservices.

More recently at Al Hilal Invest I built and documented new backend services and production REST APIs in NodeJS and TypeScript on AWS (S3/EC2), led the system-design and code reviews, and dockerized the backend to make deployments reliable and repeatable. Alongside this I have built distributed, event-driven systems in my own projects — microservices over gRPC and NATS JetStream with PostgreSQL persistence, deployed on Kubernetes — so designing scalable APIs and optimising database efficiency is familiar territory, not a stretch.

I am especially drawn to a fast-moving scale-up where I can contribute from day one on performance, scalability and integration challenges. I am based in Ilmenau, Germany, available full-time now and fully set up for remote work, and as this is a remote engagement there is no relocation or sponsorship overhead — I am Blue Card-ready, so an employer need only issue the contract. I would welcome the chance to discuss how I can help scale your platform.

\closing{Sincerely,}{Radheem Bin Razi}

\clearpage
\selectlanguage{ngerman}

%% ===== DEUTSCH =====
\recipient{XpertDirect}{Personalabteilung}{}

\opening{Sehr geehrtes Hiring-Team,}

Backend-Services zu entwickeln, die unter hohem Transaktionsvolumen schnell und korrekt bleiben, ist die Arbeit, die ich am besten beherrsche. Bei Bluefin Exchange war ich Design-Verantwortlicher für die zentralen Services einer durchsatzstarken Krypto-Börse — einer Domäne, in der jeder API-Aufruf Geld bewegt, Latenz keine Nachsicht kennt und Zuverlässigkeit nicht verhandelbar ist. Das ist dieselbe Art von Herausforderung wie hinter einer digitalen Commerce-Plattform mit hohem Volumen, die Kundinnen und Kunden in ganz Europa bedient, und genau die hier beschriebene Rolle.

Über zwei Jahre lang verantwortete ich bei Bluefin zentrale Backend-Services auf Basis von NodeJS, TypeScript, Go und PostgreSQL. Ich trieb Systemdesign-Verbesserungen eigenständig voran, prüfte Code- und Designentscheidungen des Teams und hielt Release-Prozesse und Dokumentation in gutem Zustand, während ich Junior-Engineers betreute. Die Plattform lief auf Docker und Kubernetes, wobei Kafka die ereignisgesteuerte Nachrichtenübermittlung zwischen den Services übernahm — unmittelbar relevant für Ihre Nice-to-haves rund um ereignisgesteuerte Architekturen und Microservices.

Zuletzt entwickelte und dokumentierte ich bei Al Hilal Invest neue Backend-Services und produktive REST-APIs in NodeJS und TypeScript auf AWS (S3/EC2), leitete die Systemdesign- und Code-Reviews und dockerisierte das Backend, um Deployments zuverlässig und wiederholbar zu machen. Daneben habe ich in eigenen Projekten verteilte, ereignisgesteuerte Systeme aufgebaut — Microservices über gRPC und NATS JetStream mit PostgreSQL-Persistenz, ausgerollt auf Kubernetes — sodass der Entwurf skalierbarer APIs und die Optimierung der Datenbankeffizienz vertrautes Terrain sind und keine Hürde.

Besonders reizt mich ein schnell wachsendes Scale-up, in dem ich vom ersten Tag an zu Performance-, Skalierbarkeits- und Integrationsthemen beitragen kann. Ich bin in Ilmenau, Deutschland, ansässig, ab sofort in Vollzeit verfügbar und vollständig für Remote-Arbeit eingerichtet, und da es sich um eine Remote-Tätigkeit handelt, entstehen kein Umzugs- oder Sponsoring-Aufwand — ich bin Blue-Card-bereit, sodass ein Arbeitgeber lediglich den Vertrag ausstellen muss. Ich würde mich freuen, zu besprechen, wie ich dazu beitragen kann, Ihre Plattform zu skalieren.

\closing{Mit freundlichen Grüßen,}{Radheem Bin Razi}

\end{document}
